\begin{abstract}
Cross-attention transformers and other multimodal vision-language models excel at grounding and generation; however, their extensive, full-precision backbones make it challenging to deploy them on edge devices. Memory-augmented architectures enhance the utilization of past context; however, most works rarely pair them with aggressive edge-oriented quantization. We introduce BitMar, a quantized multimodal transformer that proposes an external human-like episodic memory for effective image-text generation on hardware with limited resources. BitMar utilizes 1.58-bit encoders, one for text (BitNet-style) and one for vision (DiNOv2-based), to create compact embeddings that are combined and used to query a fixed-size key-value episodic memory. During vector retrieval, the BitNet decoder applies per‑layer conditioning, which increases the contextual relevance of generated content. The decoder also employs attention sinks with a sliding‑window mechanism to process long or streaming inputs under tight memory budgets. The combination of per-layer conditioning and sliding-window attention achieves a strong quality–speed trade–off, delivering competitive captioning and multimodal understanding at low latency with a small model footprint. These characteristics make BitMar well-suited for edge deployment.




%TOO LONG
% Cross-attention transformers and other multimodal vision-language models excel at grounding and generation. However, they typically require full-precision, large-capacity backbones, which makes them challenging to deploy on devices with limited resources. At the same time, memory-augmented architectures can facilitate the utilization of context over time. However, most systems do not combine them with aggressive quantization that is well-suited for edge settings. This paper presents BitMar, a novel quantized multimodal transformer architecture that integrates human-like episodic memory for efficient cross-modal understanding and generation on edge devices. The model comprises parallel quantized encoders for language and vision: a BitNet-style text encoder and a DiNOv2-based vision encoder, both operating with 1.58-bit weight quantization to minimize computational and memory overhead. Each encoder produces a 768-dimensional latent representation, which is fused via a lightweight cross-modal interaction module to form a unified multimodal embedding. This embedding serves as a query to a human-like episodic memory module implemented as a fixed-size key–value matrix, enabling retrieval of contextually relevant representations from prior data episodes. Retrieved memory vectors are projected into per-layer conditioning signals that modulate a BitNet decoder, allowing generation to be informed by both the current multimodal input and accumulated episodic knowledge. Furthermore, to support unbounded context lengths without growing resource usage, attention sinks are injected while maintaining a sliding window over the most recent tokens. This design achieves a balance between representational capacity and deployment efficiency, supporting competitive multimodal understanding and text generation performance while maintaining low latency and a compact model footprint suitable for edge deployment scenarios.
\end{abstract}

% \newenvironment{keyword}{\begin{center}\bfseries Keywords:Sentence Simplification \sep Generative Model \sep Large Language Model \sep Dual Latent.}{\end{center}}
\noindent\textbf{Keywords:} TinyVLM, Episodic memory, EdgeAI, Quantization.